%%%%%%%%%%%%%%%%%%%%%%%%%%%%%%%%%%%%%%%%%%%%%%%%%%%%%%%%%%%%%%%%%%%%%%%%%%%%%%%
%% Multi-Width Parametric Bitvectors (PBV_n) --- based on the OOPSLA'26 paper.
%% Shared macros and generated plot data live in
%% chapters/parametric-bv/pbv-preamble.tex, which thesis.tex inputs once before
%% the mono-width chapter and this chapter.
%%%%%%%%%%%%%%%%%%%%%%%%%%%%%%%%%%%%%%%%%%%%%%%%%%%%%%%%%%%%%%%%%%%%%%%%%%%%%%%
\chapter{Multi-Width Parametric Bitvectors}
\label{ch:multi-width}
\epigraph{Listening in wild places, we are audience to conversations in a language not our own}{Braiding Sweetgrass}

\paragraph*{Attribution.} This chapter is based on \emph{Sound and Complete
Solving for Multi-Width Parametric Bitvectors via Principled
Reductions}~\citep{bhat2026multiwidth}, conditionally accepted at OOPSLA 2026.
\emph{My contributions:} I generalized the procedure from the mono-width to the
multi-width setting, and recognized that encoding symbolic widths as bitmasks yields an
equisatisfiable reduction from $\text{PBV}_n$ to $\text{PBV}_1$. I developed the
theoretical extension, mechanized the soundness proof, and implemented the
reflection and preprocessing frontends.

\section{Overview}
\label{sec:mwbv:overview}
% Two paragraphs, essentially no citations.
% P1: the question this chapter answers, where it sits in the thesis argument,
%     and the result in one plain sentence.
% P2: the approximations made to the theoretically-justified algorithm, then
%     the empirical payoff.

\lettrine{T}{he} previous chapter decided bitvector predicates in the
\emph{mono-width} fragment $\text{PBV}_1$: a single symbolic width. That solver
came first, since it was immediately obvious to me how to build it.
Only once the mono-width procedure was in hand did it become clear to me that the far more general
\emph{multi-width} theory $\text{PBV}_n$, whose expressions may mix arbitrarily
many symbolic widths, reduces \emph{back} to it. This chapter develops that
generalization: an equisatisfiable reduction that encodes symbolic widths as
bitmasks, collapsing the many-width problem, with only linear blowup, into a
single call to the mono-width machinery of \cref{ch:mono-width}.

Fixed-width bitvector solvers (\qfbv{}) are fast and ubiquitous, but the theory
of \emph{parametric} bitvectors
(PBV), where widths are symbolic, is much less well understood in its full
generality. The theory of \emph{multi-width} PBV, where expressions may
involve $n$ distinct symbolic widths ($\text{PBV}_n$), is particularly
challenging: the only previously existing complete approach for bounded PBV
(where all widths have a concrete upper bound) is exhaustive enumeration,
requiring one \qfbv{} solver call for each of the exponentially many width
assignments. This is a significant bottleneck in tools, such as
Alive2~\citep{lopes2021alive2} and \hydra{}~\citep{hydra}, that
formally reason about compiler optimizations.

To address this problem, we first prove that any $\text{PBV}_n$ formula can be
reduced to an equisatisfiable \emph{mono-width} ($\text{PBV}_1$) formula with
only a linear increase in formula size. The key idea is to encode symbolic
widths as \emph{bitmasks}. That single reduction opens two paths---a bounded
solver that needs only one \qfbv{} call, and, by composing with the automata of
\cref{ch:mono-width}, a sound and complete decision procedure for a fragment
strictly larger than the prior state of the art---and the next section follows
both.

\section{The Algorithmic Idea: Masks and Binary Values}
\label{sec:mwbv:idea}

We give a self-contained overview of the algorithmic ideas, before diving into
the formal development in the next sections.
%
Informally, a $\text{PBV}_n$ formula is a bitvector formula with $n$ width
variables, each ranging over the naturals. Naive enumeration, i.e.,
instantiating every combination of widths up to a bound~$o$, requires on the
order of~$o^n$ solver calls, quickly becoming impractical (e.g., $64^3 =
262{,}144$ calls for three widths). Our goal is to collapse \emph{all} these
width combinations into a single bitvector universe so that one solver call
suffices. We explain the key idea via a worked example, which we will then
generalize.

Consider the following formula, which states that sign extending a
zero extended value is equivalent to zero extending it directly, when
the widths are compatible:
\[
\begin{aligned}
  &\forall (u~v~w : \N)~(a : \BV{u}),~ u < v \to v \leq w \to\\
  &\qquad \sext(\zext(a, u, v), v, w) = \zext(a, u, w)
\end{aligned}
\]
We will translate this into a single \qfbv{} query which, for a given bound~$o$
on the widths, proves the identity for all widths up to~$o$.
%
Fix a global bitwidth~$o$ (say $o = 64$); every variable from the original
formula ($u, v, w, a$) will be a new variable of type ~$\BV{o}$ in the new
formula.

The width variables $u, v, w : \N$ become bitvector variables
$u', v', w' : \BV{o}$, each storing the width value directly.
Whenever we need to constrain a value to~$u$ bits, we compute the
\emph{mask} $(1 \bitshl u') - 1 = 2^{u'} - 1$, which has exactly~$u'$
ones in the low-order bits: for instance, $u' = 3$ gives
$(1 \bitshl 3) - 1 = 7 = \langle\redbox{00}\bluebox{111}\rangle$.

Next, the bitvector variable $(a : \BV{u})$ becomes $(a' : \BV{o})$, with the
\emph{pre-masking} constraint $a' \bitand ((1 \bitshl u') - 1) = a'$ asserting
that its bits above position~$u$ are zero.
The invariant that $a$ is masked to have bits only up to $u$ active is
the key idea of the translation: it lets every subsequent operation work
uniformly in~$\BV{o}$, while still respecting the original width $u$.
%
The hypotheses translate directly: since $u'$ and $v'$ store the widths
themselves, the ordering hypothesis $u < v$ becomes $u' \bitult v'$,
and $v \leq w$ becomes $\neg(w' \bitult v')$. Signed comparisons are translated similarly,
using the most significant bit of the variables to reduce to unsigned comparisons (\cref{mwbv:fig:mono-translation}).

Next, zero extension $\zext(a, u, w)$
embeds a $u$-bit value into $w$~bits by padding with zeros. In the
translated world, every value already lies in~$(\BV{o})$, so zero-extending to
width~$w$ is simply masking: $a' \bitand ((1 \bitshl w') - 1)$. Sign extension
$\sext(a, u, v)$ is slightly more involved: it must test the most significant
bit, fill the extension region with the sign, and mask to the target width. This
is achieved by performing bit-fiddling to compute the sign fill and extension
mask on the fly, where $\msb_u(a')$ tests the most significant bit of $a'$ at position~$u$.
% \grosser{How is $msb_u$ defined?}
%
\begin{align*}
  \denote{\mathsf{sext}(a, u, v)} \;=\;\;
  &\textbf{let } s = \mathsf{ite}(\msb_u(a'),\; \bitnot 0,\; 0)
    \quad\text{\dimtext{(sign fill)}} \\[-2pt]
  &\textbf{let } \mathit{ext} = s \bitand \bitnot((1 \bitshl u') - 1)
    \quad\text{\dimtext{(bits above width~u)}} \\[-2pt]
  &(a' \bitor \mathit{ext}) \bitand ((1 \bitshl v') - 1)
\end{align*}

Applying these rules to the full formula
(\cref{mwbv:fig:worked-example}) yields a single \qfbv{}
satisfiability problem over~$\BV{o}$.
See that the translation creates a single formula that captures the behavior of the original formula for \emph{all} widths up to~$o$ simultaneously,
without any case splits on widths.
Instantiating $o = 64$ and checking unsatisfiability proves the identity for
\emph{all widths $\leq 64$} in seconds.

\begin{figure}[htbp]
\footnotesize
\resizebox{\textwidth}{!}{%
\begin{tabular}{llll}
\toprule
\textbf{\#} & \textbf{Subexpression} & \textbf{Rule} & \textbf{Translation} \\
\midrule
\rowcolor{black!10!white}
1 & width var $u$ & bitvector var & $u' : \BV{o}$ \\
2 & width var $v$ & bitvector var & $v' : \BV{o}$ \\
\rowcolor{black!10!white}
3 & width var $w$ & bitvector var & $w' : \BV{o}$ \\
  4 & $a : \BV{u}$ & BV variable & $a' : \BV{o}$, assert $a' \bitand ((1 \bitshl u') - 1) = a'$ \\
\midrule
\rowcolor{black!10!white}
5 & $u < v$ & wlt & assert $u' \bitult v'$ \\
6 & $v \leq w$ & $\neg\,$wlt & assert $\neg(w' \bitult v')$ \\
\midrule
\rowcolor{black!10!white}
7 & $\zext(a, u, v)$ & zext & $a' \bitand ((1 \bitshl v') - 1)$ \\
8 & $\sext((7), v, w)$ & sext & $\textbf{let } s = \mathsf{ite}(\msb_v((7)),\; \bitnot 0,\; 0)$ \quad\dimtext{(sign-extend fill)} \\
  &                  &      & $\textbf{let } \mathit{ext} = s \bitand \bitnot((1 \bitshl v') - 1)$ \quad\dimtext{(bits above width $v$)} \\
  &                  &      & $((7) \bitor \mathit{ext}) \bitand ((1 \bitshl w') - 1)$ \\
\rowcolor{black!10!white}
9 & $\zext(a, u, w)$ & zext & $a' \bitand ((1 \bitshl w') - 1)$ \\
\midrule
10 & $(8) = (9)$ (negated) & equality & $(8) \neq (9)$ \\
\bottomrule
\end{tabular}
}
\caption{Worked example: translating
$\forall (u\;v\;w : \N)\;(a : \BV{u}),\;
u < v \to v \leq w \to \sext(\zext(a, u, v), v, w) = \zext(a, u, w)$
from $\text{PBV}_n$ to a single \qfbv{} query over $\BV{o}$,
previewing the rules formalized in \cref{mwbv:fig:mono-translation}.
Primed variables ($u', v', w', a'$) denote the translated bitvector counterparts.
Unsatisfiability at $o = 64$ proves the identity for all widths $\leq 64$.}
\label{mwbv:fig:worked-example}
\end{figure}

\paragraph{The General Principle.}
%
The worked example illustrates a general pattern: every $\text{PBV}_n$
operation maps to a fixed sequence of $\BV{o}$ operations, with widths stored
directly as bitvector variables and masks computed on the fly as
$(1 \bitshl w') - 1$.
The pre-masking invariant—each value is masked to its logical width—ensures that
the translation is homomorphic: no case splits on widths are needed, and width
comparisons reduce to ordinary unsigned bitvector comparisons. The formal
development (Figure~\ref{mwbv:fig:mono-translation},
page~\pageref{mwbv:fig:mono-translation}) makes this precise; the primed notation is
dropped there, as the context is unambiguous.

This translation yields two paths, depending on whether we wish to prove the
formula up to a bound, or to completely prove it when it lies in the
linear-bitwise fragment.

\paragraph{Path 1: Bounded Model Checking via a Single \qfbv{} Call.}
%
After reducing $\text{PBV}_n$ to $\text{PBV}_1$, we instantiate the unique
width parameter~$o$ to a concrete value (e.g., $64$), yielding a standard
\qfbv{} formula that any off-the-shelf SMT solver can discharge in a single
call.  This handles the full theory, including multiplication and
concatenation.  The cost drops from~$o^k$ solver calls under naive enumeration
to exactly one call.

\paragraph{Path 2: Sound and Complete Decision via Automata.}
%
For the linear-bitwise fragment, we can do better: prove validity for \emph{all}
widths simultaneously.  The idea is to extend bitvectors to infinite bitstreams
that are computable by finite state transducers, and decide the resulting
formula by model checking. One nuance is that widths are directly represented by
their associated masks, which can be seen as their unary encodings.

\medskip
Overall, the same idea of reducing widths to masks provides two
reductions, one to a single \qfbv{} query for bounded model checking, and one to
a complete automata-theoretic decision procedure for the unbounded case. We will
now formalize the syntax and semantics of the multi-width theory and then state
these reductions precisely.


\section{A Unified Presentation of Zero, One, Many Width Parametric Bitvector Theory}
\label{sec:mwbv:theory}

We now define syntax and semantics for a family of multi-width
parametric bitvector theories $\text{PBV}_n$, where~$n$ is the number of
bitvector width variables.
%
If we restrict width expressions to have no width variables ($\text{PBV}_0$), we
recover the standard theory of fixed-width bitvectors (\qfbv{}). If we allow
only one width variable, we recover the mono-width parametric bitvector formula from the previous chapter.
%
Thus, this unifies the presentation of all three theories, and allows us to
state our reductions: an equisatisfiable translation from $\text{PBV}_n$ to
$\text{PBV}_1$, followed by bounded model checking via a single \qfbv{}
call, all in a unified framework, parameterized by $n$.

\subsection{Syntax}
\label{mwbv:sec:syntax}

Our theory is a multi-sorted language (\cref{mwbv:fig:syntax}, left) with two sorts of variables:
\emph{width variables}~$x_w$ (ranging over natural numbers) and
\emph{bitvector variables}~$x$ (ranging over bitvectors whose widths are given by width expressions).
%
Width expressions~$w$ comprise constants, variables, and simple operations
such as min, max, and addition.
%
Bitvector expressions contain the standard operations from the SMT-LIB 2 \cite{barrett2010smt}
format, such as arithmetic operations, bitwise operations, append,
and extension operations. For example $\mathsf{zext}(a, w_1, w_2)$ extends the
width of the bitvector~$a$ from~$w_1$ to~$w_2$ by adding zeros, and
$\mathsf{sext}(a, w_1, w_2)$ by duplicating the sign bit of~$a$.
%
Finally, the predicates we consider are Boolean combinations of equality and
signed and unsigned order predicates.
For bounded solving, our algorithm allows all operations from SMT-LIB. For unbounded solving,
our algorithm restricts to operations within our fragment.

\begin{figure}[htbp]
\begin{minipage}{0.55\textwidth}
{\footnotesize
\resizebox{\textwidth}{!}{%
\begin{tabular}{llll}
%\toprule
% \multicolumn{4}{@{}l}{\textit{Width Expressions}} \\[2pt]
$w$ & $::=$ & $\mathsf{wconst}(c)$ & constant width $c \in \N$ \\
      & $\mid$ & $\Var(x_w)$ & width variable \\
      & $\mid$ & $\mathsf{wmax}(w_1, w_2)$ & maximum of two widths \\
      & $\mid$ & $\mathsf{wmin}(w, w_2)$ & minimum of two widths \\
      & $\mid$ & $\mathsf{wadd}(w_1, w_2)$ & width addition \\[4pt]
%\multicolumn{4}{@{}l}{\textit{Bitvector Expressions}} \\[2pt]
$a, b$ & $::=$ & $\Var(x, w)$ & bitvector variable of width $w$ \\
    & $\mid$ & $\mathsf{Const}(c, w)$ & constant $c$ at width $w$ \\
    & $\mid$ & $\mathsf{add}(a, b, w)$ & addition at width $w$ \\
    & $\mid$ & $\mathsf{mul}(a, b, w)$ & multiplication at width $w$ \\
    & $\mid$ & $\mathsf{nand}(a, b, w)$ & bitwise NAND at width $w$ \\
    & $\mid$ & $\mathsf{append}(a, w_1,  b, w_2)$ & append two bitvectors \\
    & $\mid$ & $\zext(a, w_1, w_2)$ & zero-extend from $w_1$ to $w_2$ \\
    & $\mid$ & $\sext(a, w_1, w_2)$ & sign-extend from $w_1$ to $w_2$ \\[4pt]
%\multicolumn{4}{@{}l}{\textit{Predicate Expressions}} \\[2pt]
$p$ & $::=$ & $\mathsf{eq}(a, b, w)$ & bitvector equality \\
    & $\mid$ & $\mathsf{ult}(a, b, w)$ & unsigned less-than \\
    & $\mid$ & $\mathsf{slt}(a, b, w)$ & signed less-than \\
    & $\mid$ & $\mathsf{and}(p, q)$ & conjunction \\
    & $\mid$ & $\mathsf{or}(p, q)$ & disjunction \\
    & $\mid$ & $\mathsf{weq}(u, v)$ & width equality \\
    & $\mid$ & $\mathsf{wlt}(u, v)$ & width less-than \\
%\bottomrule
\end{tabular}
}
}
\end{minipage}%
%\hfill
\begin{minipage}{0.45\textwidth}
{\footnotesize
\begin{tabular}{ll}
$\Delta \defn$ \dimtext{(Width Typing Contexts)} \\
$\Gamma \defn$ \dimtext{(Bitvector Typing Contexts)} \\
$\Delta \vdash w$ \dimtext{($w$ well-typed in $\Delta$)} \\
$\Delta; \Gamma \vdash x : \BV{w}$ \dimtext{($x$ has width $w$ in $\Gamma$)} \\
$\Delta \vdash \IsEq(v, w)$ \dimtext{($v$, $w$ equal in $\Delta$)}
\end{tabular}
% \vspace{0.5em}
\begin{tabular}{@{}l@{}}
\rowcolor{black!10!white}
\inferrule{ }{\Delta \vdash \mathsf{wconst}(c) : \N} \\[1em]
\inferrule{ x_w \in \Delta}{\Delta \vdash x_w : \N} \\[1em]
\rowcolor{black!10!white}
\inferrule{ \Delta \vdash w_1 : \N \quad \Delta \vdash w_2 : \N}{\Delta \vdash \mathsf{wmax}(w_1, w_2) : \N} \\[1em]
\inferrule{\Delta \vdash w }{\Delta; \Gamma \vdash \mathsf{Const}(c, w) : \BV{w}} \\[1em]
\rowcolor{black!10!white}
\inferrule{\Delta \vdash w \quad (x, w) \in \Gamma}{\Delta; \Gamma \vdash x : \BV{w}} \\[1em]
\inferrule{\Delta; \Gamma \vdash x : \BV{w_x} \\ \Delta; \Gamma \vdash y : \BV{w_y} \\ \Delta \vdash \IsEq(w_x, w_y)}{\Delta; \Gamma \vdash x + y : \BV{w_x}} \\[1em]
\rowcolor{black!10!white}
\inferrule{\Delta; \Gamma \vdash x : \BV{w_1} \quad \Delta \vdash w_2  }{\Delta; \Gamma \vdash \zext(x, w_1, w_2) : \BV{w_2}} \\
\end{tabular}
}
\end{minipage}
\caption{Left: Syntax fragment of the multi-width parametric bitvector theory.
Right: Representative typing rules. Note the stratified typing contexts for widths ($\Delta$) and bitvector variables ($\Gamma$).}
\label{mwbv:fig:syntax}
\label{mwbv:fig:typing}
\end{figure}


\subsection{Type System}
\label{mwbv:sec:type-system}

Bitvectors in our theory are dependently typed: the type $\BV{w}$ is
parameterized by a width value~$w$ of type~$\N$.
%
Formally, we stratify our typing contexts into two levels (\cref{mwbv:fig:typing},
right).
%
We first have a context~$\Delta$ for width variables, of type~$\N$.
%
We also have a context~$\Gamma$ for bitvector variables, which depend on the
width variables. This context maps each bitvector variable to a well-typed width
expression. We consider typing contexts to be unordered sets of variables, and
environments as functions from typing contexts to values.
% \footnote{In our implementation,
% we use ordered lists (instead of unordered sets) for better definitional equality,
% but the theoretical type system obeys all structural rules.}
As there are expressions in types, we need a notion of equality for these
expressions for type checking:
Consider the rule for addition, where we require that the two bitvector operands
have the same width, in the sense that they denote the same natural numbers
under every instantiation of the width variables, denoted by $\IsEq(w_1, w_2)$.
%
In our implementation, this is approximated by a syntactic equality check.

\subsection{Semantics}
\label{mwbv:sec:semantics}

\begin{figure}[htbp]
\centering
\small
\setlength{\tabcolsep}{3pt}
\renewcommand{\arraystretch}{1.2}
\begin{tabular}{@{}p{12.5cm}@{}}
\toprule
\textbf{Semantics} (relative to a model~$\rho$) \\
\midrule
\textit{Width expressions}: $\rho(w) \in \N$ \\[2pt]
\rowcolor{black!10!white}
$\rho(\mathsf{wconst}(c)) \defn c$ \\
$\rho(\Var(w)) \defn \rho(w)$ \\
\rowcolor{black!10!white}
$\rho(\mathsf{wmax}(k, l)) \defn \max(\rho(k), \rho(l))$ \\
$\rho(\mathsf{wmin}(k, l)) \defn \min(\rho(k), \rho(l))$ \\
\rowcolor{black!10!white}
$\rho(\mathsf{wadd}(k, c)) \defn \rho(k) + c$ \\[4pt]
\midrule
\textit{Bitvector expressions}: $\rho(e) \in \BV{\rho(w)}$ for $e : \BV{w}$ \\[2pt]
\rowcolor{black!10!white}
$\rho(\Var(x, w)) \defn \rho(x)$ \\
$\rho(\mathsf{Const}(c, w)) \defn c \bmod 2^{\rho(w)}$ \\
\rowcolor{black!10!white}
$\rho(\mathsf{add}(a, b, w)) \defn (\rho(a) + \rho(b)) \bmod 2^{\rho(w)}$ \\
$\rho(\mathsf{nand}(a, b, w)) \defn \lnot(\rho(a) \bitand \rho(b)) \bmod 2^{\rho(w)}$ \\
\rowcolor{black!10!white}
$\rho(\zext(a, v, w)) \defn$ zero-extend $\rho(a)$ from width $\rho(v)$ to width $\rho(w)$ \\
$\rho(\sext(a, v, w)) \defn$ sign-extend $\rho(a)$ from width $\rho(v)$ to width $\rho(w)$ \\[4pt]
\midrule
\textit{Predicates}: the satisfaction relation $\rho \vDash p$ \\[2pt]
\rowcolor{black!10!white}
$\rho \vDash \mathsf{eq}(a, b, w) \iff \rho(a) = \rho(b)$ \\
$\rho \vDash \mathsf{ult}(a, b, w) \iff \natofbv{\rho(a)} < \natofbv{\rho(b)}$ \\
\rowcolor{black!10!white}
$\rho \vDash \mathsf{slt}(a, b, w) \iff \intofbv{\rho(a)} < \intofbv{\rho(b)}$ \\
$\rho \vDash \mathsf{and}(p, q) \iff \rho \vDash p \ \text{and}\ \rho \vDash q$ \\
\rowcolor{black!10!white}
$\rho \vDash \mathsf{or}(p, q) \iff \rho \vDash p \ \text{or}\ \rho \vDash q$ \\
$\rho \vDash \mathsf{weq}(u, v) \iff \rho(u) = \rho(v)$ \\
\rowcolor{black!10!white}
$\rho \vDash \mathsf{wlt}(u, v) \iff \rho(u) < \rho(v)$ \\
\bottomrule
\end{tabular}
\caption{Semantics of the multi-width parametric bitvector theory, relative to a
model~$\rho$ that maps each width variable to a natural number and each bitvector
variable to a bitvector of the corresponding width. Bitvector and width
expressions are interpreted by the function~$\rho(\cdot)$; predicates by the
satisfaction relation~$\rho \vDash \cdot$.}
\label{mwbv:fig:semantics}
\end{figure}

We now give standard semantics for our theory, collected in
\cref{mwbv:fig:semantics}.
% A key difference between the semantics we adopt versus model-theoretic semantics is that
We only assign semantics to well-typed terms,
which lets us avoid having to define
partial functions for operations such as bitvector addition for bitvectors of
two different widths. Our models~$\rho$ assign a natural number $\rho(x_w)$ for
every width variable~$x_w$. This function is extended to associate a natural
number $\rho(w)$ to every width expression~$w$ in the expected way.

A model~$\rho$ also associates a bitvector $\rho(x)$ of width~$\rho(w)$ for a
bitvector $x : \BV{w}$.
%
We extend this to a function $\rho$ over well-typed bitvector expressions using
the standard semantics of bitvector operators such as addition or extension. We
obtain a function which associates a bitvector expression $a$ of type~$\BV{w}$
with a bitvector~$\rho(a) \in \BV{w}$.

Finally, we define when a model satisfies a formula, which we denote with $\rho
\vDash p$ in the expected way: for instance, $\rho \vDash \mathsf{eq}(a, b, w)$ when
$\rho(a) = \rho(b)$.
%
We write $\vDash p$ when all models satisfy~$p$.
The full definition is given in \cref{mwbv:fig:semantics}, and its mechanized
counterpart in the accompanying Lean development.

\section{Equisatisfiable Reduction: Multi-Width to Mono-Width}
\label{sec:mwbv:reduction}

We now formalize the reduction previewed in \cref{sec:mwbv:idea}:
an equisatisfiable translation from $\text{PBV}_n$ to $\text{PBV}_1$,
translating multi-width $\text{PBV}_n$ problems (as defined in \cref{sec:mwbv:theory})
into a mono-width $\text{PBV}_1$ problem.
The second reduction, bounded model checking via a single \qfbv{} call, follows immediately: fixing a concrete width $o$ proves
the property for all widths $\leq o$.
Composing the two, the entire multi-width formula maps to
one \qfbv{} satisfiability check, so any off-the-shelf SMT solver can discharge it,
with no external looping or enumeration to make multiple \qfbv{} calls, which is a significant improvement over the state of the art.

The core idea is to represent each width~$w$ as a bitvector variable storing the
width value directly, and to compute \emph{masks} on the fly whenever a
masking operation is needed:
the mask corresponding to a width~$w$ is $(1 \bitshl w) - 1 = 2^{w} - 1$, a
bitvector with~$w$ ones in the lower bits.
%
For example, a width variable holding $u = 3$ yields the mask
$(1 \bitshl 3) - 1 = 7 = \langle \redbox{00}\bluebox{111} \rangle$,
with ones at the three low-order positions~$0$--$2$.

\subsection{Syntax-Directed Translation to Mono-Width PBV}

\begin{figure}[htbp]
\small
\setlength{\tabcolsep}{3pt}
\renewcommand{\arraystretch}{1.1}
\centering
\begin{tabular}{@{}p{4.5cm}p{8cm}@{}}
\toprule
\textbf{Multi-width PBV} & \textbf{Mono-width PBV} $\denote{\cdot}$ \\
\midrule
\multicolumn{2}{@{}l}{\textit{Variable Declarations} (preconditions asserted)} \\[2pt]
\rowcolor{black!10!white}
$\textsf{Var}(x_w) : \N$     & $(x_w : \BV{o})$ \\
$\textsf{Var}(x, w) : \BV{e}$           & $(x : \BV{o})$, \quad assert $(x \bitand \denote{e}^\mask = x)$ \\[4pt]
\midrule
\multicolumn{2}{@{}l}{\textit{Width Expressions} (widths stored directly as bitvectors)} \\[2pt]
\rowcolor{black!10!white}
  $\textsf{Var}(x_w)$          & $x_w$ \\
$\textsf{const}(c)$           & $c$ \\
\rowcolor{black!10!white}
$\textsf{min}(u, v)$           & $\mathsf{ite}(\denote{u} \bitult \denote{v},\; \denote{u},\; \denote{v})$ \\
$\textsf{wadd}(u, v)$           & $\denote{u} + \denote{v}$ \\[4pt]
\midrule
\multicolumn{2}{@{}l}{\textit{Bitvector Expressions} (mask result to width via $\denote{w}^\mask = (1 \bitshl \denote{w}) - 1$)} \\[2pt]
\rowcolor{black!10!white}
$\textsf{Var}(a, w)$       & $a$ \quad\dimtext{(pre-masked by variable declaration)} \\
$\textsf{Const}(x, w)$     & $x \bitand \denote{w}^\mask$ \\
\rowcolor{black!10!white}
$\textsf{add}(a, b, w)$    & $(\denote{a} + \denote{b}) \bitand \denote{w}^\mask$ \\
$\textsf{nand}(a, b, w)$   & $\bitnot(\denote{a} \bitand \denote{b}) \bitand \denote{w}^\mask$ \\
\rowcolor{black!10!white}
$\textsf{zext}(a, v, w)$       & $\denote{a} \bitand \denote{w}^\mask$ \\
$\textsf{sext}(a, v, w)$       & $\textbf{let } s = \mathsf{ite}(\msb_v(\denote{a}),\; \bitnot 0,\; 0)$ \quad\dimtext{(sign-extend fill)} \\
                       & $\textbf{let } \mathit{ext} = s \bitand \bitnot \denote{v}^\mask$ \quad\dimtext{(bits above width $v$)} \\
                       & $(\denote{a} \bitor \mathit{ext}) \bitand \denote{w}^\mask$ \\
\rowcolor{black!10!white}
$\textsf{append}(a, v, b, w)$  & $(\denote{a} \bitshl \denote{w}) \bitor \denote{b}$ \\
% \rowcolor{black!10!white}
                       % & \dimtext{result width $v{+}w$ encoded as $\denote{v} \cdot \denote{w}$} \\[4pt]
\midrule
\multicolumn{2}{@{}l}{\textit{Predicate Expressions}} \\[2pt]
\rowcolor{black!10!white}
$\textsf{eq}(a, b, w)$  & $\denote{a} = \denote{b}$ \quad\dimtext{(inputs pre-masked)} \\
$\textsf{ult}(a, b, w)$  & $\denote{a} \bitult \denote{b}$ \quad\dimtext{(inputs pre-masked)} \\
\rowcolor{black!10!white}
% if the signs are equal,
% then the result is that of unsigned comparison of 'a < b', flipped if a (and b) are negative.
% otherwise the signs are different, the negative number is less than the positive number, so return if 'a' is negative.
  $\textsf{slt}(a, b, w)$ & $\mathsf{ite}(\msb_w(\denote{a}) = \msb_w(\denote{b})$, \\
\rowcolor{black!10!white}
                          & \quad $\msb_w(\denote{a})\oplus\textsf{ult}(\denote{a}, \denote{b})$, \\
\rowcolor{black!10!white}
                          & \quad $\msb_w(\denote{a}))$ \\
$\textsf{and}(p, q)$     & $\denote{p} \land \denote{q}$ \\
\rowcolor{black!10!white}
$\textsf{or}(p, q)$      & $\denote{p} \lor \denote{q}$ \\
$\textsf{weq}(u, v)$                & $\denote{u} = \denote{v}$ \\
\rowcolor{black!10!white}
$\textsf{wlt}(u, v)$                & $\denote{u} \bitult \denote{v}$ \\[2pt]
\bottomrule
  \multicolumn{2}{@{}l}{\dimtext{where $\msb_w(a) \defn (a \bitand ((1 \bitshl \denote{w}) \gg 1)) \neq 0$.}} \\
\end{tabular}
\caption{Syntax-directed translation from $\text{PBV}_n$ to $\text{PBV}_1$.}
\label{mwbv:fig:mono-translation}
\end{figure}

% \paragraph{Width Operations via Bit Manipulation}
% Width operations on masks translate to standard bitvector operations.
% Given masks $m, n$ (each of the form $w^{\pot} - 1$), we have
% $\min(m, n) = m \bitand n$,
% $\max(m, n) = m \bitor n$, and
% \mbox{$m + 1 = (m \bitshl 1) \bitor 1$}.
% Width predicates are equally direct:
% $m$ and $n$ are equal iff $m \bitxor n = 0$,
% and $m < n$ iff $\bitnot n \bitand m = 0$.

We combine these observations into the syntax-directed translation
of \cref{mwbv:fig:mono-translation}, which covers every operation
from our multi-width calculus (\cref{mwbv:fig:syntax}).
Given a multi-width formula over width variables $x_w^1, \ldots, x_w^k$ and bitvector variables
$\textsf{Var}(x_1, w_1), \ldots, \textsf{Var}(x_n, w_n)$,
we consider a global width $o$ and translate all variables and expressions to $\BV{o}$.
Width variables are translated to bitvector variables $w : \BV{o}$ storing the
width value directly (we denote the translated width with the same name).
%
Width expressions~$w$ are interpreted as bitvectors $\denote{w}$, and the
corresponding mask $\denote{w}^\mask$ is computed on the fly as
$(1 \bitshl \denote{w}) - 1 = 2^{\denote{w}} - 1$.

Bitvector variables $\textsf{Var}(x, w)$ become $x : \BV{o}$, constrained to be
correctly masked to their logical width: $x \bitand \denote{w}^\mask = x$.
%
Bitvector expressions are translated by masking their result back to the result
width. Predicates and Boolean connectives translate homomorphically.


\subsection{Equisatisfiability}\label{mwbv:sec:bounded-proof}

We prove that the translation $\denote{\cdot}$ of \cref{mwbv:fig:mono-translation} is correct:
every $\text{PBV}_n$ formula~$p$ is equisatisfiable with its
translation~$\denote{p}$. We prove this through a single relation between the
models of the two formulas. Intuitively, a mono-width model represents a
multi-width one by storing each narrow bitvector as its zero-extension to the
global width~$o$, keeping the logical bits unchanged and forcing every bit above
the width to zero. We call such paired models \emph{compatible}
(\cref{mwbv:def:compat}), show that compatible models satisfy exactly the same
formulas (\cref{mwbv:lem:correspondence}), and conclude equisatisfiability by
building a compatible partner from multi-width to mono-width and vice versa (\cref{mwbv:thm:equisat}).

\begin{definition}[Compatibility]\label{mwbv:def:compat}
Let $\rho$ be a multi-width model of~$p$, and let $W$ be the
  maximum width of any subexpression in $p$. Formally, we define $W \defn \max\{\rho(w) : w
\text{ a width subexpression of } p\}$.
  Let $\rhob$ be a mono-width model over the translated global width~$o$. We call
$\rho$ and $\rhob$ \emph{compatible}, written $\rho \approx \rhob$, when $o > W$
and
\begin{align*}
  \rhob(x_w) \;&=\; \ofNat_o(\rho(x_w)) && \text{for every width variable } w, \\
  \rhob(x)   \;&=\; \mathsf{zext}_o(\rho(x))     && \text{for every bitvector variable } x : \BV{w}.
\end{align*}
Here $\ofNat_o(n)$ writes the natural number~$n$ as a $\BV{o}$, and
$\mathsf{zext}_o$ zero-extends a bitvector of width $w$ to a width $o$.
The latter is exactly masking as zero-extension is implemented by a bitmask:
$\mathsf{zext}_o(\rho(x))$ equals $\rhob(x) \bitand \denote{w}^\mask$.
\end{definition}

The constraint $o > W$ leaves room for every width and every mask: each width
value $\rho(w) \le W < o$ is representable in~$\BV{o}$, and each mask $(1 \bitshl
\denote{w}) - 1$ is computed without overflow.

\begin{lemma}[Correspondence]\label{mwbv:lem:correspondence}
If $\rho \approx \rhob$ then
\begin{enumerate}
  \item for every width expression~$w$,\quad $\rhob(\denote{w}) = \ofNat_o(\rho(w))$;
  \item for every bitvector expression~$a : \BV{w}$,\quad $\rhob(\denote{a}) = \mathsf{zext}_o(\rho(a))$;
  \item for every formula~$q$,\quad $\rho \vDash q \iff \rhob \vDash \denote{q}$.
\end{enumerate}
\end{lemma}
\begin{proof}
All of the claims are proven by an induction on the structure of the well-typed expression,
reading the translation from \cref{mwbv:fig:mono-translation}.

\emph{(1) Widths.} Width variables hold by compatibility, and
$\textsf{const}(c)$ is immediate. For $\textsf{wadd}$ and $\textsf{min}$, the
bound $o > W$ keeps every result below~$2^o$, so bitvector addition and
$\bitult$ agree with addition and comparison on~$\N$.

\emph{(2) Bitvectors.} The key fact is that masking by $\denote{w}^\mask$ at the
wide width~$o$ commutes with performing the operation natively at the logical
width~$\rho(w)$. Addition illustrates the pattern: assuming the claim for $a$
and~$b$,
\noindent
\begin{align*}
  \rhob(\denote{\textsf{add}(a,b,w)}) &= (\mathsf{zext}_o(\rho(a)) + \mathsf{zext}_o(\rho(b))) \bitand \denote{w}^\mask \\
  &= \mathsf{zext}_o\big((\rho(a) + \rho(b)) \bitand (2^{\rho(w)} - 1)\big) \\
  &= \mathsf{zext}_o(\rho(\textsf{add}(a,b,w))),
\end{align*}
as the mask erases exactly the high bits, where the overflow from the larger width $o$
can disagree with the smaller width $\rho(w)$. The remaining cases (bitwise, extensions, append)
follow the same argument, where masking the wide result reproduces the narrow result.

\emph{(3) Formulas.} By parts~(1) and~(2), each atomic predicate compares the
same values in both models. For example, $\rho \vDash \textsf{eq}(a,b,w)$ iff
$\rho(a) = \rho(b)$ iff $\mathsf{zext}_o(\rho(a)) = \mathsf{zext}_o(\rho(b))$ iff
$\rhob \vDash \denote{\textsf{eq}(a,b,w)}$. Boolean connectives
remain unchanged, and thus the truth value of the entire formula is preserved.
\end{proof}

\begin{theorem}[Equisatisfiability]\label{mwbv:thm:equisat}
A $\text{PBV}_n$ formula~$p$ is satisfiable iff its translation $\denote{p}$ is satisfiable.
\end{theorem}
\begin{proof}
We will use the correspondence lemma (\cref{mwbv:lem:correspondence})
to pair the multi-width formula $p$ with a compatible model of the mono-width formula $\denote{p}$ (and vice versa).

\emph{(multi-width $\Rightarrow$ mono-width)} Let $\rho \vDash p$. Set $o = W + 1$ and define $\rhob$ by
$\rhob(x_w) = \ofNat_o(\rho(x_w))$ and $\rhob(x) = \mathsf{zext}_o(\rho(x))$.
Then $\rho \approx \rhob$ by construction, and every masking precondition
$x \bitand \denote{w}^\mask = x$ holds because $\rhob(x)$ is a zero-extension.
The lemma gives $\rhob \vDash \denote{p}$.

\emph{(multi-width $\Leftarrow$ mono-width)} Let $\rhob \vDash \denote{p}$ over global width~$o$. Define
$\rho$ for widths by interpreting the bitvectors $\rhob(w) : \BV{o}$ as natural numbers: $\rho(w) = \toNat(\rhob(w))$.
Similarly, translate the bitvectors $x : \BV{o}$ as bitvectors of width
$\rho(w)$ by masking: $\rho(x) = \ofNat_w (\toNat (\rhob(x)))$.
The precondition $\rhob(x) \bitand \denote{w}^\mask = \rhob(x)$ ensures that $\toNat \rho(x) = \toNat \rhob(x)$, so
this truncation does not lose information (as \toNat~is injective), and thus $\rho \approx \rhob$. The lemma gives $\rho \vDash p$.
\end{proof}

\begin{remark}[Choice of the global width $o$]\label{mwbv:rem:width-choice}
The two maps in the proof of \cref{mwbv:thm:equisat} are a one-sided inverse.
Starting from a multi-width model~$\rho$, building $\rhob$ and then truncating recovers $\rho$ exactly.
So, this translation is lossless.

On the other hand, the global width~$o$ is not fully determined by~$p$, as any $o$
that is sufficient to not cause overflow ($o \ge W + 1$) yields a valid compatible model.
Our construction simply takes the smallest such width, $o = W + 1$, for efficiency.
\end{remark}


\section{Extending Automata-Based PBV Solving to Multiple Widths}
\label{sec:mwbv:automata}
The previous section reduced $\text{PBV}_n$ to $\text{PBV}_1$ via bounded model
checking, yielding a single \qfbv{} call that proves validity for all widths up
to a bound~$o$.
%
For an \emph{unbounded} decision procedure that proves validity for
\emph{all} widths simultaneously, we turn to automata.
%
As the full theory is undecidable, we need to restrict ourselves to a fragment
of~$\text{PBV}_n$ which contains linear arithmetic operations as well as bitwise
operations.
%
We combine the idea of the reduction from $\text{PBV}_n$ to $\text{PBV}_1$ with
the verified mono-width solver of~\citet{bhat2025certified} to build a decision
procedure which supports multiple widths and zero and sign extension.


\subsection{The linear-bitwise fragment of $\text{PBV}_n$}

The theory we consider is the fragment of $\text{PBV}_n$ which restricts the
bitvector operations to linear arithmetic (addition, subtraction,
multiplication by a constant), bitwise operations (and, or, xor, etc),
size-changing operations (truncation, sign and zero extension), and comparisons (equality, unsigned and signed less-than).

\subsection{From Bounded Masks to Automata-Encodable Masks}
\label{mwbv:sec:mask-substitution}

The translation of \cref{sec:mwbv:reduction} computes width masks on the fly as
$\denote{w}^\mask = (1 \bitshl \denote{w}) - 1$, a left-shift by a variable
amount (\cref{mwbv:fig:mono-translation}).  Variable shifts fall outside the
linear-bitwise fragment, so the translation cannot be fed directly to an existing
automata-based mono-width solver~\citep{bhat2025certified}.
%
We exploit a classical bit-manipulation identity:
%
\begin{lemma}\label{mwbv:lem:mask-char}
  A bitvector~$m$ equals $2^k - 1$ for some $k \geq 0$
  if and only if $m \bitand (m + 1) = 0$.
\end{lemma}
%
\noindent
For each width variable~$w$ in the original $\text{PBV}_n$ formula, we
can introduce a fresh bitvector variable~$m_w$ and assert
$m_w \bitand (m_w + 1) = 0$.
We can then perform the translation of \cref{mwbv:fig:mono-translation}
with every occurrence of~$\denote{w}^\mask$ replaced by~$m_w$.
The resulting formula is mono-width (all bitvectors live in~$\BV{o}$)
and lies entirely within the linear-bitwise fragment,
since every operation is either linear arithmetic, bitwise, or a constant shift.
%
This demonstrates that automata are expressible enough to decide linear-bitwise
formulas with multiple widths.
%
Instead of using the above encoding and generating a mono-width formula, which
would take significant effort to prove, we instead use this idea to directly
translate multi-width formulas to automata by extending the existing
infrastructure~\citep{bhat2025certified}.
%
We also conjecture that giving custom encodings of zero and sign extension yields
smaller automata, which are easier to model check.

%% \subsection{Background: The Automata Theoretic Approach to Mono-Width PBV}
%% 
%% We describe a recent decision procedure~\citep{bhat2025certified} for the
%% mono-width linear-fragment of PBV (i.e., without the three size-changing
%% operations listed above).
%% 
%% \begin{wrapfigure}[5]{r}{0.4\textwidth}
%% {
%% % \footnotesize
%% \vspace{-2em}
%% \setlength{\tabcolsep}{4pt}
%% \begin{tabular}{r|cccccc}
%%   \rowcolor{black!10!white}
%%           & $\dots$ & $b_4$ & $b_3$ & $b_2$ & $b_1$ & $b_0$ \\
%%   $x$           & $\dots$ & 0 & 0 & 0 & 1 & 1 \\
%%   \rowcolor{black!10!white}
%%   $y$           & $\dots$ & 0 & 0 & 0 & 0 & 1 \\
%%   \rowcolor{black!5!white}
%%   $x \bitand y$ & $\dots$ & 0 & 0 & 0 & 0 & 1 \\
%%   \rowcolor{black!10!white}
%%   $y \bitand x$ & $\dots$ & 0 & 0 & 0 & 0 & 1 \\
%% \end{tabular}
%% }
%% \end{wrapfigure}
%% 
%% For an arbitrary width $w$, suppose we have the bitvector predicate $P_w \defn
%% \forall (x~y: \BV{w}), x \bitand y = y \bitand x$.
%% Our objective is to prove that the predicate $P_\ast \defn \forall w, P_w$ holds.
%% The first reduction step is to see that we can instead prove the equivalent predicate
%% $\Phi \defn \forall (x~y : \Bitstream), x \bitand y = y \bitand x$, which is
%% over infinite bitstreams, and where bitwise operations are defined in the
%% expected way.
%% %
%% This first reduction converts an equality over parametric bitvectors into an equality over bitstreams.
%% %
%% We imagine these bitstreams as infinite sequences of bits indexed from $0$ to $\infty$,
%% laid out on a strip. For a particular instantiation, suppose $x = 3\#5 = \langle 00011\rangle$
%% (i.e. bitvector of value $3$, bitwidth $5$) and $y = 1\#5 = \langle 00001\rangle$.
%% These become infinite bitstreams by extending them with an infinite sequence of $0$s.
%% %
%% Next, to decide $\Phi$, we need to be able to have a finite representation for
%% the bitstreams $x \bitand y$ and $y \bitand x$.
%% %
%% We compute these bitstreams via finite state transducers, on which many
%% properties are decidable.
%% %
%% For illustration, we depict the bitstreams for $x \bitand y$ and $y \bitand x$
%% on the right.
%% 
%% Now, for the predicate $x \bitand y = y \bitand x$, we consider equality itself as a bitstream.
%% Here, the equality predicate is a stream of $T$ and $F$ values, where $T$ indicates that
%% the two input bitstreams are equal up to that position,
%% and $F$ indicates that the two input bitstreams have differed at some position.
%% For our running example, since the expression is true, the equality bitstream is all $T$s.
%% Now, the question of whether $\Phi$ holds reduces to whether the bitstream for $(x \bitand y) = (y \bitand x)$ is always $T$.
%% 
%% \begin{figure}[tbp]
%% %\begin{minipage}[t]{0.3\textwidth}
%% \begin{tabular}{lll}
%%   \begin{tikzpicture}[mwbvautomaton, font=\footnotesize,node distance=2cm, on grid, auto, background rectangle/.style={fill=black!10!white}, show background rectangle, rounded corners=0.3cm]
%% 
%%     % Left automaton: (x, y) -> x
%%       \node[state, initial, accepting] (x0) {$x_0$};
%%       \path[->]
%%         (x0) edge [loop above] node {$[0,0]/0, [0,1]/0$} ()
%%              edge [loop below] node {$[1,0]/1, [1,1]/1$} ();
%%       \node[above=1.9cm of x0, font=\small] {$(x, y) \mapsto x$};
%% \end{tikzpicture}
%% &
%%   \begin{tikzpicture}[mwbvautomaton, font=\footnotesize,node distance=2cm, on grid, auto, background rectangle/.style={fill=black!10!white}, show background rectangle, rounded corners=0.3cm]
%%       \node[state, initial, accepting] (y0) {$y_0$};
%%       \path[->]
%%         (y0) edge [loop above] node {$[0,0]/0, [1,0]/0$} ()
%%              edge [loop below] node {$[0,1]/1, [1,1]/1$} ();
%%       \node[above=1.9cm of y0, font=\small] {$(x, y) \mapsto y$};
%% \end{tikzpicture}
%% &
%% %\end{minipage}
%% %\begin{minipage}[t]{0.3\textwidth}
%%   \begin{tikzpicture}[mwbvautomaton, font=\footnotesize,node distance=2cm, on grid, auto, background rectangle/.style={fill=black!10!white}, show background rectangle, rounded corners=0.3cm]
%%     % \backgroundrectangle[fill=red!10!white]{} % optional extra settings
%%     \node[state, initial, accepting] (x0) {$x_0$};
%%     \path[->]
%%         (x0) edge [loop above] node {$[0,0]/0, [0,1]/0, [1, 0]/0$} ()
%%              edge [loop below] node {$[1,1]/1$} ();
%%     \node[above=1.9cm of x0, font=\small] {$(p, q) \mapsto p \bitand q$};
%% \end{tikzpicture}
%% \end{tabular}
%% \caption{Automata for $x$, $y$, and $x \bitand y$, operations with no state needed.}
%% \end{figure}


% \paragraph{Finite State Transducers for Mono-Width Parametric Bitvectors}
% 
% The next step is to decide $\Phi$.
% For this, we model the input bitstreams $x$ and $y$, as well as all intermediate expressions
% (e.g., $x \bitand y$) as the output of finite state transducers.
% These transducers receive a tuple of bits with one tuple element for each variable, at each timestep,
% and produce the bitstream corresponding to the expression's value at that timestep.
% As an example, the expression $x$ receives a tuple of bits as input $[x_i, y_i]$, and produces output $x_i$.
% Similarly, the automaton for the subexpression $y$ receives a tuple of bits $[x_i, y_i]$ as input and produces $y_i$ as output.
% 
% For the subexpressions $x \bitand y$ and $y \bitand x$,
% we first see that there is a finite state transducer that computes $(p, q)\mapsto p \bitand q$.
% We obtain the automata for $x \bitand y$ as well as $y \bitand x$ by composing the automata for $x$, $y$ and the $\bitand$ transducer.
% While in this case, the automata appear identical, in general, they will differ based on the input expressions.
% We can construct an automaton for the predicate $\Phi$ by composing the automata for $x \bitand y$ and $y \bitand x$.
% 
% 
% 
% \begin{figure}[tbp]
% \begin{tabular}{ll}
% \begin{tikzpicture}[mwbvautomaton, font=\footnotesize, shorten >=1pt, node distance=2cm, on grid, auto, background rectangle/.style={fill=black!10!white}, show background rectangle, rounded corners=0.3cm]
%     % \backgroundrectangle[fill=blue!10!white]{} % optional extra settings
%    % States
%    \node[state, initial] (eq) {=};
%    \node[state, right=of eq] (neq) {≠};
% 
%    % Loops
%    \path[->]
%      (eq) edge[loop above] node {[1,1]/T} ()
%           edge[loop below] node {[0,0]/T} ()
%      (neq) edge[loop above] node {[$\star$, $\star$]/F} ();
% 
%    % Transitions
%    \path[->]
%      (eq) edge[bend left] node {[0,1]/F} (neq)
%           edge[bend right] node[below] {[1,0]/F} (neq);
%     \node[above=1.9cm of eq, font=\small] {$(p, q) \mapsto p = q$};
% \end{tikzpicture} &
% \begin{tikzpicture}[mwbvautomaton, font=\footnotesize, shorten >=1pt, node distance=2cm, on grid, auto, background rectangle/.style={fill=black!10!white}, show background rectangle, rounded corners=0.3cm]
%    % States
%    \node[state, initial, accepting] (eq) {=};
%    \node[state, right=of eq] (neq) {$\neq$};
% 
%    % Loops on = state
%    \path[->]
%     (eq) edge[loop above] node{$[0, 0]/T, [0, 1]/T$} ()
%          edge[loop below] node{$x:[1, 0]/T, [1, 1]/T$} ();
% 
%    % Loop on != state
%    \path[->]
%     (neq) edge[loop above] node{$[\star, \star] / F$} ();
% 
%    % Transition from = to !=
%    \path[->]
%     (eq) edge[bend left, above] node{} (neq);
%     \node[above=1.9cm of eq, font=\small] {$(x, y) \mapsto (x \bitand y) = (y \bitand x)$};
% \end{tikzpicture}
% \end{tabular}
% \caption{Automata for equality and the final predicate $\Phi$.}
% \end{figure}
% 
% 
% So far, all the automata we have constructed have been memoryless. The next automaton that we construct for $(p, q) \mapsto p = q$
% is more interesting, since it uses state to keep track of whether the two input bitstreams have differed at any prior timestep.
% The idea is that this automaton starts in the $(=)$ state, and remains in this state as long as the two input bitstreams are equal.
% If, at any timestep, the two input bitstreams differ, the automaton transitions to the $(\neq)$ state,
% and remains in this state forever. This matches our intuition of equality:
% two bitstreams that are disequal at any position must remain disequal forevermore.
% %
% Finally, we compose the automata for $x \bitand y$ and $y \bitand x$ with the automaton for equality,
% giving us the final automaton for the predicate $\Phi$. See that in the final automaton,
% the $F$ state is unreachable, so $\Phi$ holds.
% 
% \subsection{Extending to Multiple Widths: New Automata}

Recall that ater the masking translation for a given input formula, (\cref{mwbv:sec:mask-substitution}), all width variables
become mask bitstreams.
The mono-width automata handle arithmetic, bitwise operations, and masking,
but we still need automata for the operations unique to the multi-width setting.
We illustrate the construction with zero extension.
Using the running example of
\[
\begin{aligned}
  &\forall (u~v~w : \N)~(a : \BV{u}),~ u < v \to v \leq w \to\\
  &\qquad \sext(\zext(a, u, v), v, w) = \zext(a, u, w)
\end{aligned}
\]

We encode the width variables $u, v, w$ as mask bitstreams $m_u, m_v, m_w$.
These bitstreams are such that the first $u$ bits of $m_u$ are $1$s and the rest are $0$s, and similarly for $m_v$ and $m_w$.
Note that $m_u = 2^u - 1$, and similarly for $m_v, m_w$; each is the mask for its corresponding width.

\begin{figure}[htbp]
\begin{tabular}{ll}
\begin{tikzpicture}[mwbvautomaton, font=\footnotesize, shorten >=1pt, node distance=2cm, on grid, auto, background rectangle/.style={fill=black!10!white}, show background rectangle, rounded corners=0.3cm]
    % \backgroundrectangle[fill=blue!10!white]{} % optional extra settings
   % States
   \node[state, initial] (T) {$\geq$};
   \node[state, right=of T] (F) {F};

   % Loops
   \path[->]
     (T) edge[loop above] node {[1,1]/T, [1, 0]/T} ()
          edge[loop below] node {[0,0]/T} ()
     (F) edge[loop above] node {[$\star$, $\star$]/F} ();

   % Transitions
   \path[->]
     (T) edge[bend left] node {[0,1]/F} (F);
   \node[above=1.9cm of T, font=\small] {$(z, w) \mapsto z \geq w$};
\end{tikzpicture}\hspace{3em}
&
\begin{tikzpicture}[mwbvautomaton, font=\footnotesize, shorten >=1pt, node distance=2cm, on grid, auto, background rectangle/.style={fill=black!10!white}, show background rectangle, rounded corners=0.3cm]
    % \backgroundrectangle[fill=blue!10!white]{} % optional extra settings
   % States
   \node[state, initial] (T) {$\zext$};
   \node[state, right=of T] (F) {0};

   % Loops
   \path[->]
    (T) edge[loop above] node {[1,1,1]/1, [0,1,1]/0} ()
          edge[loop below, black!10!white] node {\phantom{[0,0]/T}} () % hack: invisible arrow to make the two figures the same height
    (F) edge[loop above] node {[$\star$, $\star$, $\star$]/0} ();

   % Transitions
   \path[->]
    (T) edge[bend left, above] node {[$\star$,0, $\star$]/F} (F)
    (T) edge[bend right, below] node {[$\star$,$\star$,0]/F} (F);
   \node[above=1.9cm of T, font=\small] {$(x, v, w) \mapsto \zext(x, v, w)$};
\end{tikzpicture}
\end{tabular}
\caption{Automata for multi-width mask operators,  $z \geq w$ and $\zext(x, v, w)$.
Note the mixture of mask bitstreams (for $z, w, v$) and binary bitstreams (for $x$).
}
\end{figure}

Now, width comparisons such as $u < v$ and $v \leq w$ can be encoded as comparisons of the corresponding mask bitstreams $m_u < m_v$ and $m_v \leq m_w$.
This works because the map $k \mapsto 2^k - 1$ is monotone, so $u < v$ if and only if $m_u < m_v$.

Next, for zero extension, we construct an automaton for the operation $\zext(x, m_v, m_w)$,
which takes as input a bitstream $x$ and two mask bitstreams $m_v$ and $m_w$,
and produces the zero extension of $x$ from width $v$ to width $w$.
The key idea is that zero extension simply masks the input bitstream $x$ with $m_v$ to keep the bits within width $v$,
and then masks this result with $m_w$ to get the final zero-extended result.
So, the zero extension $\zext(x, m_v, m_w)$ can be computed as $x \bitand m_v \bitand m_w$,
which is in the linear-bitwise fragment and can be computed by a finite state transducer.
The other size-changing operation, sign extension, needs a \msb{} automaton.
This can be computed by using both the original mask $m_w$, and the mask shifted by one, $m_{w-1}$,
to mask away just the top bit. Then, we check whether this masked bitvector is nonzero to determine the msb.
Formally, $\msb(x)$ can be computed as $(x \bitand m_w \bitand \bitnot  m_{w-1} \neq 0)$.
Sign extension can be computed in terms of msb and other automata-expressible operations.
As explained, these extend the prior chapter by also requiring width masks to be available as bitstreams,
which allows one to encode zero and sign extension into the linear-bitwise fragment, which is automata decidable.

\begin{figure}[tbp]
\newcommand{\addCarry}{\textsf{addCarry}}
\newcommand{\addOut}{\textsf{addOut}}
\newcommand{\subBorrow}{\textsf{subBorrow}}
\small
\setlength{\tabcolsep}{3pt}
\renewcommand{\arraystretch}{1.1}
\centering
\resizebox{\textwidth}{!}{%
\begin{tabular}{@{}p{3.8cm}p{10.8cm}@{}}
\toprule
\textbf{Syntax} & \textbf{Semantics} $\denote{\cdot}[n]$ \\
\midrule
\multicolumn{2}{@{}l}{\textit{Width Expressions} \newstuff{(mask encoding)}} \\[2pt]
\rowcolor{black!10!white}
$\newstuff{\textsf{const}(c)}$, $\newstuff{\textsf{Var}(v)}$
& $\newstuff{\denote{\textsf{const}(c)}[n] \defn [n < c]}$ \\
\rowcolor{black!10!white}
& $\newstuff{\denote{\textsf{Var}(v)}_\omega[n] \defn \omega(v)[n]}$ \\[2pt]
$\newstuff{\textsf{min}(v, w)}$
& $\newstuff{\denote{\textsf{min}(v, w)}_\omega[n] \defn \denote{v}_\omega[n] \bitand \denote{w}_\omega[n]}$ \\[4pt]
\midrule
\multicolumn{2}{@{}l}{\textit{Bitvector Expressions} (binary encoding, \newstuff{width-indexed})} \\[2pt]
\rowcolor{black!10!white}
$\textsf{Var}(a, \newstuff{w})$, $\textsf{Const}(x, \newstuff{w})$
& $\denote{\textsf{Var}(a, \newstuff{w})}_{\newstuff{\omega}, \beta}[n] \defn \beta(a)[n] \newstuff{\bitand \denote{w}_\omega[n]}$ \\
\rowcolor{black!10!white}
& $\denote{\textsf{Const}(x, \newstuff{w})}_{\newstuff{\omega}, \beta}[n] \defn \textsf{testbit}(x, n) \newstuff{\bitand \denote{w}_\omega[n]}$ \\[2pt]
$\newstuff{\textsf{zext}(a, v, w)}$
& $\newstuff{\denote{\textsf{zext}(a, v, w)}_{\omega, \beta}[n] \defn \denote{a}_{\omega, \beta}[n] \bitand \denote{\textsf{min}(v, w)}_\omega[n]}$ \\[2pt]
\rowcolor{black!10!white}
$\textsf{add}(a, b, \newstuff{w})$
& $\denote{\textsf{add}(a, b, \newstuff{w})}_{\omega, \beta}[n] \defn \addOut(\denote{a}, \denote{b}, 0)[n] \newstuff{\bitand \denote{w}_\omega[n]}$ \\[2pt]
$\textsf{nand}(a, b, \newstuff{w})$
& $\denote{\textsf{nand}(a, b, \newstuff{w})}_{\omega, \beta}[n] \defn \lnot(\denote{a}[n] \bitand \denote{b}[n]) \newstuff{\bitand \denote{w}_\omega[n]}$ \\[4pt]
\midrule
\multicolumn{2}{@{}l}{\textit{Predicate Expressions} (true $=$ stream of 1s)} \\[2pt]
\rowcolor{black!10!white}
$\textsf{eq}(a, b, \newstuff{w})$
& $\denote{\textsf{eq}(a, b, \newstuff{w})}_{\newstuff{\omega}, \beta}[n] \defn \newstuff{\forall i \leq n,\, (\denote{w}_\omega[i] {=} 1) \to} \denote{a}[i] {=} \denote{b}[i]$ \\[2pt]
$\textsf{ult}(a, b, \newstuff{w})$
& $\denote{\textsf{ult}(a, b, \newstuff{w})}[n] \defn \subBorrow(\denote{a}, \denote{b})[n] \newstuff{\bitand \denote{w}_\omega[n]}$ \\[2pt]
\rowcolor{black!10!white}
$\textsf{and}(p, q)$, $\textsf{or}(p, q)$
& $\denote{\textsf{and}(p, q)}[n] \defn \denote{p}[n] \land \denote{q}[n]$ \quad (similarly for $\textsf{or}$) \\[2pt]
$\newstuff{\textsf{weq}(v, w)}$, $\newstuff{\textsf{wlt}(v, w)}$
& $\newstuff{\denote{\textsf{weq}(v, w)}_\omega[n] \defn \forall i \leq n,\, \denote{v}_\omega[i] = \denote{w}_\omega[i]}$ \\
\rowcolor{black!10!white}
& $\newstuff{\denote{\textsf{wlt}(v, w)}_\omega[n] \defn \exists i \leq n,\, \denote{v}[i] {>} \denote{w}[i] \land (\forall j < i,\, \denote{v}[j] {=} \denote{w}[j])}$ \\[4pt]
\midrule
\multicolumn{2}{@{}>{\raggedright\arraybackslash}p{\dimexpr 3.8cm+10.8cm+6pt\relax}@{}}{\textit{Auxiliary}: $\addOut(x, y, c)[i] \defn x[i] \oplus y[i] \oplus \addCarry(x, y, c)[i{-}1]$, where $\addCarry$ computes the carry bit.} \\
\multicolumn{2}{@{}>{\raggedright\arraybackslash}p{\dimexpr 3.8cm+10.8cm+6pt\relax}@{}}{\phantom{\textit{Auxiliary}:} $\subBorrow(x, y)[i] \defn \addCarry(x, \lnot y, 1)[i]$ (borrow from subtraction $x - y$).} \\
\bottomrule
\end{tabular}
}
\caption{Encoding of bitvector operations as regular languages.
Widths are masks, values are binary, predicates are bitstreams.
Extensions are \newstuff{highlighted}.}
\label{mwbv:fig:encodings}
\end{figure}


\subsection{From Transducers to Model Checking}

The key to decidability is to encode the bitstreams as outputs of finite state transducers.
Encoding widths as mask bitstreams allows us to perform width-based bit-masking on the bitvector operations,
% thereby extending the mono-width encodings of prior work~\cite{bhat2025certified} to the multi-width setting.
Our extensions to the original calculus from the previous chapter are shown in \cref{mwbv:fig:encodings}, with extensions \newstuff{highlighted}.

Given a finite state transducer $f$ which encodes the denotation of a bitvector predicate $\phi$,
we can decide the truth value of $\phi$ by ensuring that all reachable states of $f$ produce an output of $1$.
We use a mechanized $k$-induction solver to check this property of $f$.
As $k$-induction is sound and our encodings are sound, our overall decision procedure is sound.
We reuse the $k$-induction solver developed in the previous chapter.
We also have an implementation of a solver based on IC3, which is not proof producing, but is much more efficient on certain problem structures.


\subsection{Soundness \& Completeness of Automata Encoding}

The decision procedure is sound and complete:
\begin{theorem}
  Let $p$ be a predicate in the linear-bitwise fragment of $\text{PBV}_n$.
  %
  Then $\vDash p$ if and only if the associated transducer produces a stream
  of~$T$.
\end{theorem}
%
\noindent
This result is mechanized using Lean. The idea of the proof
is similar to that of the previous chapter. We reuse the library we have built for expressing finite state machines,
and the $k$-induction solver with proof certificates.
We build the new syntax encodings, translations to automata, and soundness and completeness proofs for the translation.
The proof takes the new width constraints into account: we relate expressions with their associated
transducers, by stating that for a given environment, a bitvector term~$a$
produces the bitstream (by completing it with zeros) as its associated
transducer.
%
We can then lift this result to predicates, and prove that a predicate~$p$ is
true in an environment iff the associated transducer only produces~$T$ with the
environment as input.

\section{Solver Implementation and Mechanization}
\label{sec:mwbv:implementation}

We implement our decision procedures in Lean~4, producing
both a tactic (\ourtactic{}) for interactive theorem proving and a standalone CLI solver.

\begin{figure}[htbp]
    \centering
    % https://docs.google.com/drawings/d/1OXxLRow6Swbr4vzk9T75vbfxAL1mwaDQp9q4tDeWonY/edit?usp=sharing
    \includegraphics[height=0.42\textheight]{chapters/multi-width-bv/images/Meta-Concrete-Picture-Multi-Width-BV-Vertical.pdf}
    \caption{End-to-end tactic pipeline.
    We reflect the goal state, and coordinate type-erased (at meta-level) \& intrinsically well-typed (at object-level) syntax.}
    \label{mwbv:fig:pipeline}
\end{figure}

The implementation has two layers of syntax representation,
two solver backends and two frontends (\cref{mwbv:fig:pipeline}).

\subsection{Intrinsically Well-Typed Syntax}

Our implementation uses two syntax representations, connected by a type-erasing projection.
The first is an \emph{intrinsically well-typed} syntax (implemented as
\texttt{Term} and \texttt{Predicate})
whose types are indexed by width and term contexts,
ensuring well-formedness by construction;
for instance, a binary operation can only be formed
when both subterms share the same width expression.
This representation is used for the soundness proof,
which proceeds by induction on well-typed terms.
The second is a \emph{type-erased} syntax
(implemented as \texttt{Nondep.Term} and \texttt{Nondep.WidthExpr})
that carries no dependent-type indices;
it is used at the metaprogramming level
(where manipulating dependent types is extremely painful)
and serves as the shared input to all solver backends.
A projection%
\footnote{The function \texttt{ofDepTerm} in our codebase.}
maps well-typed syntax to untyped syntax,
and the FSM construction for well-typed terms is \emph{defined}
by composing this projection with the FSM construction on untyped terms:
$\texttt{Predicate} \rightarrow \texttt{Nondep.Predicate} \rightarrow
\texttt{FSM}$. This ensures both representations produce identical automata,
which is critical for soundness: proof certificates obtained from the untyped
FSM are directly applicable to the well-typed term.

\subsection{Verified Unbounded Solver By Automata Reduction}

The unbounded solver targets the linear-bitwise fragment decidable by
model-checking of transducers (\cref{sec:mwbv:theory}). It proceeds in
two steps.

First, the solver builds a finite-state transducer%
\footnote{Via the function \texttt{mkTermFsmNondep}.}
from the untyped syntax tree.
The correctness theorem states that for every well-typed term,
the FSM output matches the bitstream semantics.
This is proved by induction on the well-typed syntax,
using a reusable reasoning principle about FSM outputs
in terms of inductive invariants over FSM states.

Second, the solver decides whether the constructed FSM is safe.
Here we offer two backend choices.
The verified option invokes a $k$-induction solver with cycle breaking built on
previous work by \citet{bhat2025certified}:
the solver calls \cadical{} and produces LRAT proof certificates,
which are checked by Lean's verified certificate checker
(part of \texttt{bv\_decide}~\cite{leanbv}),
yielding a fully verified end-to-end proof.
The new unverified option invokes the rIC3 model checker~\cite{ric3} as an
external process; this is typically faster than $k$-induction, and the
architecture is designed to be able to verify the certificates generated by
rIC3, which contain witnesses for the invariants.

\subsection{Bounded Solver By $\text{PBV}_n$ To $\text{PBV}_1$ Reduction}

The bounded solver handles the full $\text{PBV}_n$ language, including
operations outside the automata-encodable fragment. Unlike the unbounded solver,
which offers a choice of backends, the bounded solver is a fixed three-stage
pipeline, though it can be implemented as a frontend to any solver which
supports the \qfbv{} theory.
%
The first stage applies the equisatisfiable reduction
from $\text{PBV}_n$ to $\text{PBV}_1$ (\cref{sec:mwbv:reduction}).%
\footnote{Implemented as \texttt{toSingleWidthNondepTerm}.}
The second stage translates mono-width syntax to the input format
for \texttt{bv\_decide}.%
\footnote{Via \texttt{toBVLogicalExpr}, producing Lean's \texttt{BVLogicalExpr} type.}
This pipeline is unverified.
The third stage hands the resulting \qfbv{} problem to \texttt{bv\_decide},
Lean's verified bitblaster~\cite{leanbv}.

The trust placed in this pipeline therefore rests on
(1)~correctness of the unverified reduction, and
(2)~the bounded model checking assumption that the chosen bound suffices.
  In certain cases, this assumption might well be true, where the upper bound is known a priori,
  such as when reasoning about fixed-width machine integers.
The key advantage of this technique is that it handles the full PBV language,
and requires only a \emph{single} solver call,
as opposed to techniques that enumerate width combinations
(exponential in the number of widths).

\subsection{Tactic Frontend For Lean \& Command Line Executable For SMT-LIB}

We implement two frontends: a Lean tactic for interactive theorem proving
and a standalone executable for SMT-LIB input.

\noindent\textbf{Lean tactic (unbounded solver).}
When the user invokes \ourtactic{} on a goal that lies in the linear-bitwise
fragment of $\text{PBV}_n$, the tactic traverses the goal's AST at the meta
level, building a type-erased syntax tree together with typing contexts. This
avoids the pain of manipulating dependent types during metaprogramming. The
well-typed term is then constructed at the object level from the untyped syntax
and the separately tracked typing information. The tactic invokes the unbounded
solver on the untyped syntax tree and obtains a proof certificate (an LRAT
string added to the Lean environment as a constant).
% A ``master theorem'' then closes the goal by combining three components:
% (1)~the definitional equality of the reflected syntax with the goal,
% (2)~the equivalence of FSM semantics with bitstream semantics, and
% (3)~the certificate checker returning \texttt{true}.
% Each definitional equality is taken as a propositional hypothesis (proved by \texttt{refl})
% rather than used directly, producing cleaner error messages
% when reflection fails.

\noindent\textbf{SMT-LIB executable.}
We also create a standalone CLI that consumes SMT-LIB~2 files
in a parametric bitvector fragment, with configurable choice of backend (bounded or unbounded).

\subsection{Preprocessing \& SMT-LIB Interoperability}

Our theory is built on well-typed expressions whose widths are natural numbers
(\cref{mwbv:sec:syntax}), whereas the SMT-LIB PBV
proposal~\cite{bit-precise-parametric-bv} uses integer-valued width expressions.
To bridge the gap, our preprocessor performs three transformations:
(1)~SMT-LIB's $\texttt{pzero\_extend}(d, x)$, which extends by a delta~$d$,
is rewritten as $\texttt{zext}(x, w, w+d)$ in our target-width convention;
(2)~each width subtraction $u - v$ is replaced by a fresh width variable~$k$
with the constraint $k + v = u$; and
(3)~bitvector if-then-else is eliminated into implications.

\noindent\textbf{Natural-number width encoding.}
Extension operations in our formalization take the \emph{target width}
as argument---mirroring the convention used in Lean and LLVM---so that every
argument is non-negative by construction.
Concretely, SMT-LIB's $\texttt{pzero\_extend}(d, x)$ becomes
$\texttt{zext}(x, w, w{+}d)$, and similarly for sign-extend.

\noindent\textbf{If-then-else elimination.}
Our automata construction operates on a linear fragment and does not
directly support bitvector if-then-else.
The preprocessor eliminates each $\texttt{bvIte}(c,\, t_1,\, t_2)$
by introducing a fresh bitvector variable~$r$ of the same width
and adding the implications
$c \Rightarrow r = t_1$ and $\lnot c \Rightarrow r = t_2$
as conjunctive preconditions; the original if-then-else is then replaced by~$r$.
This is an equisatisfiable transformation that brings the formula
into the if-then-else-free fragment handled by our solvers.

\noindent\textbf{Width-subtraction elimination.}
The subtraction-elimination pass replaces each width subtraction $a - b$
with a fresh width variable~$(d : \N)$ and adds the precondition $d + b = a$,
effectively encoding subtraction that truncates to zero, since we are subtracting natural numbers (monus).
Identical subtractions are deduplicated, keeping the variable count small.

We do not perform any more preprocessing such as normalization up to associativity-commutativity of bitvector expressions,
as we wish to evaluate the solving technique directly, without transformations that may normalize the problem and thereby make it trivial to solve.

\section{Evaluation}
\label{sec:mwbv:evaluation}

\begin{figure}[htbp!]
\newcommand{\myspace}{1.5em}
\centering
\vspace{\myspace}
\begin{subfigure}[t]{\textwidth}
\centering
\includegraphics[width=\textwidth]{chapters/multi-width-bv/plots/output/alive-cactus/cactus.pdf}
\caption{PBV-Alive (\AliveCactusNumProblemsTotal{} mono-width problems).
BMC-Ours saturates the dataset; both unbounded solvers outperform CVC5-berger.}
\label{mwbv:fig:eval-all-cactus}
\end{subfigure}

\vspace{\myspace}

\begin{subfigure}[t]{\textwidth}
\centering
\includegraphics[width=\textwidth]{chapters/multi-width-bv/plots/output/sext-zext-instcombine/cactus.pdf}
\caption{InstCombine (\SextZextInstcombineNumProblemsTotal{} multi-width problems).
BMC-Ours saturates the dataset despite multiple widths; BMC-Naive collapses from enumeration blowup.}
\label{mwbv:fig:eval-sext-zext}
\end{subfigure}

\vspace{\myspace}

\begin{subfigure}[t]{\textwidth}
\centering
\includegraphics[width=\textwidth]{chapters/multi-width-bv/plots/output/rover-cactus/cactus.pdf}
\caption{ROVER (\RoverCactusNumProblemsTotal{} problems, up to 12 symbolic widths).
Naive enumeration is completely impractical; BMC-Ours and the unbounded solvers remain effective.}
\label{mwbv:fig:eval-all-rover}
\end{subfigure}
\caption{Elapsed time against problems solved across all three datasets.
BMC-Ours saturates nearly every benchmark by using a single \qfbv{} query,
while BMC-Naive's naive enumeration collapses on multi-width problems.
Both our unbounded solvers consistently outperform CVC5-berger, proving strictly
more problems with parametric bitwidths.}
\label{mwbv:fig:eval-compiler}
\end{figure}

Two questions drive the evaluation. The first is whether the equisatisfiable
reduction from $\text{PBV}_n$ to $\text{PBV}_1$ really does eliminate the
exponential blowup of naive enumeration---whether a single \qfbv{} query suffices
to saturate standard PBV benchmarks (\cref{mwbv:sec:bounded-eval}). The second is
whether the automata-theoretic procedures for unbounded bitwidth outperform the
state of the art on parametric bitvector problems
(\cref{mwbv:sec:unbounded-eval}). We take them up in turn, on benchmarks drawn
from compiler verification and hardware optimization.

For the first, we evaluate the bounded solver on the full PBV language against a
naive enumeration-based approach. We saturate our benchmarks with a single
\qfbv{} query, while naive enumeration collapses on multi-width
problems---BMC-Ours solves
\EverybenchNumSolvedMbmc{}/\EverybenchNumProblemsTotal{} benchmarks versus
\EverybenchNumSolvedNbmc{} for BMC-Naive. For the second, we evaluate the
unbounded solvers on the automata-encodable fragment of PBV against the state of
the art CVC5-based parametric-width solver
of~\citet{bit-precise-parametric-bv}, and improve significantly on the number of
width-generic theorems solved across all three datasets, proving
\EverybenchNumSolvedMr{}/\EverybenchNumProblemsTotal{} benchmarks compared to
\EverybenchNumSolvedC{} for CVC5-berger.


\subsection{Experimental Setup}
Our experiments were performed on a system equipped with an \MwSystemSpecsProcessorName{} clocked at 4.3GHz and \MwSystemSpecsMemoryGb{}~GiB of RAM.
All experiments were executed under Ubuntu 24.04, with per-instance resource limits of \MwConfigTimeoutSec{}s of wall-clock time and \MwConfigMemoutMb{}~MB of memory.
We run \MwConfigNproc{} experiments in parallel and report wall-clock time elapsed for each solver.

\subsection{Datasets: Mono and Multi-Width Compiler and Many-Width Hardware Benchmarks}
% \grosser{Clearly setup the datasets. What was our objective in picking the data-sets,
% how do the datasets we picked satisfy the objective. For example, are we interested in compiler verification
% and hardware optimization, and if so, why? Are these BV problems representative?
% What about the benchmarks from the CVC5 people? Did they get dropped?}
Answering both questions requires benchmark datasets that are
(1)~representative of real-world applications of parametric bitvector solving,
(2)~as well-established as possible,
and (3)~challenging for existing solvers. Our datasets focus on two application domains:
compiler verification and hardware optimization. Compiler verification is a
well-known application domain for parametric bitvector solving with an
established mono-width benchmark suite (\textbf{PBV-Alive}) which we complement
with a newly-created multi-width dataset (\textbf{LLVM Instcombine (zext/sext)}). Hardware
optimization adds a set of particularly challenging many-width problems (\textbf{ROVER}).

\noindent\textbf{PBV-Alive.}
The Alive dataset from~\citet{bit-precise-parametric-bv} contains 200 mono-width parametric bitvector problems.
This is representative of the use of PBV in compiler verification, where one is interested in proving the correctness of optimizations for all bitwidths.
These are drawn from the Alive2 translation validation tool for LLVM,
covering five parts of InstCombine: \texttt{AddSub} (19), \texttt{AndOrXor} (130), \texttt{InstCombineShift} (11),
\texttt{Select} (27), and \texttt{muldivrem} (13).
Each problem has a mono-width parameter \texttt{k},
and the bitvector operations include arithmetic, bitwise operations, and comparisons.

\noindent\textbf{LLVM InstCombine (zext/sext).}
To provide compiler verification benchmarks that involve multiple widths, we create a new dataset from LLVM InstCombine patterns that involve zero and sign extensions.
We extract \SextZextInstcombineNumProblemsTotal{} multi-width problems
from the LLVM InstCombine source code for patterns that involve zero and sign extensions.
We create the dataset by parsing InstCombine test files, and running LLVM's InstCombine pass on them to receive a before/after pair of LLVM IR.
We take 20 representative transformations and manually convert them to SMT-LIB formulas, generalizing them to arbitrary width statements in the process.
We leave the original LLVM problem as a comment in the SMT-LIB file.
Then, using this as a few-shot prompt~\cite{brown2020language} for an agentic harness,\footnote{Claude Code, Opus 4.6, high effort}
we group the remaining test files into representative transformations, and have the agent convert the before/after LLVM IR into SMT-LIB formulas.
We instruct the agent to focus on problems that (a) have two or more zero/sign extend operations, and
(b) have clean generalizations to parametric widths. During the translation, we allow tool use for the agent
to verify the SMT-LIB translation by running a naive enumerative solver with a bound of $8$.
This gives us \SextZextInstcombineNumProblemsTotal{} problems from 1500 tests.

\noindent\textbf{ROVER.}
We take all 19 problems from the ROVER hardware optimization tool~\cite{DBLP:journals/tcad/CowardDC24}'s benchmark suite,
which involves up to 12 symbolic width variables with constraints, and a variety of bitvector operations.
These problems arise from RTL optimization, where ROVER uses e-graphs to perform rewriting-based equivalence checking.
These rewrites not only involve many width variables, but also have constraints that relate the widths of different bitvectors (e.g., $q \geq t$, $u \geq t$).
The large number of widths makes them a particularly demanding test on both counts: the enumeration blowup is high, and unbounded solving must scale to many width variables.

\subsection{Solvers}

To evaluate our solvers, we use the state-of-the-art algorithm for parametric bitvector solving from~\citet{bit-precise-parametric-bv} as a baseline for unbounded solving.
For bounded solving, we implement a naive enumeration-based solver as the baseline, which is the approach taken by Alive2, LLVM's translation validation tool.
We run our unbounded solver in two configurations: one using a verified $k$-induction backend, and the other using the unverified but more efficient rIC3 model checker.
BMC-Ours and BMC-Naive isolate the effect of the reduction, by comparing it against naive enumeration.
The unbounded solvers and CVC5-berger compare our automata-based approach against the state of the art.

\noindent\textbf{BMC-Ours (bounded).}
Applies our $\text{PBV}_n \to \text{PBV}_1$ reduction, then instantiates the
mono-width parameter to a concrete bound, encoding the entire problem as a
\emph{single} \qfbv{} query. We use \bvdecide{}, Lean's verified \qfbv{} solver,
since it provides a path toward an end-to-end verified decision procedure,
and is the standard solver used in the Lean community for bitvector problems.

\noindent\textbf{BMC-Naive (bounded).}
The baseline bounded model checker that enumerates all width assignments up to a given bound,
issuing one \qfbv{} query per assignment, as done by Alive2 for example.
For $k$ width variables and bound $o$, this requires $O(o^k)$ solver calls.
We use the same \bvdecide{} solver for each query, to isolate the impact of the reduction from the underlying solver performance.
This ensures that improvement in query complexity is the main driver of performance gains.

\noindent\textbf{Unbounded-Ours (ric3~$+$~k-induction).}
Both unbounded solvers share the same frontend: they construct the multi-width automaton
encoding of the $\text{PBV}_n$ formula (verified in Lean) and reduce validity to an
automaton safety property. They differ only in the model-checking backend.
The \emph{ric3} variant uses rIC3, a state of the art bit-level model checker~\cite{ric3};
it is fast but currently unverified, though its architecture supports future verification
via inductive-invariant certificates~\cite{hwmcc_certs}.
The \emph{k-induction} variant uses a $k$-induction checker mechanized in Lean,
producing fully verified proofs when it succeeds, at the cost of higher runtime.

\noindent\textbf{CVC5-berger (baseline).}
The parametric bitvector solver from~\citet{bit-precise-parametric-bv},
implemented in a fork of CVC5. We run the binary for the solver\footnote{\texttt{pbvsolver --pbv --cvc5:"nl-ext-tplanes full-saturate-quant"}}  from their artifact
\cite{zohar_2025_15143242}, which serves as the state of the art baseline for $\text{PBV}_n$ solving.

% \input{chapters/multi-width-bv/plots/output/bmc-cactus/cactus.tex}
% \subsection{Order of Magnitude Faster Bounded Model Checking for PBV}
%
% \begin{figure}[htb]
% \includegraphics[width=\textwidth]{chapters/multi-width-bv/plots/output/bmc-cactus/cactus.pdf}
% \label{mwbv:fig:eval-bmc-cactus}
% \end{figure}
%


\begin{table}[htb]
\centering
% \grosser{What does 'Unknown' mean? Maybe explain in legend or call it out-of-fragment? Also, is there a chance we can extend this and/or model these unknowns as uninterpreted functions?}
% \sid{Yes, we can eventually model these as uninterpreted functions, but I wanted to keep the TCB of our solver small.
% Every extension has the potential to introduce bugs.}
\caption{Summary of solver results across all three datasets.
\textbf{Bold} marks the best value per column within each dataset.
BMC-Ours saturates nearly every benchmark with a single \qfbv{} query, eliminating the enumeration bottleneck.
Our unbounded solvers consistently outperform CVC5-berger, proving strictly more problems for all bitwidths.
    }
\label{mwbv:tab:eval-summary}
\resizebox{\textwidth}{!}{%
\begin{tabular}{ll rrrrrr}
\toprule
Dataset & Solver & Total & Solved & Timeout & Memout & Unknown & Geomean \\
\midrule
\multirow{5}{*}{PBV-Alive}
  & BMC-Ours          & \AliveCactusNumProblemsTotal{} & \textbf{\AliveCactusNumSolvedMbmc{}} & \AliveCactusNumTimeoutMbmc{}          & \textbf{\AliveCactusNumMemoutMbmc{}} & \textbf{\AliveCactusNumUnknownMbmc{}} & \AliveCactusGeomeanMsMbmc{} \\
  & BMC-Naive         & \AliveCactusNumProblemsTotal{} & \AliveCactusNumSolvedNbmc{}          & \AliveCactusNumTimeoutNbmc{}          & \AliveCactusNumMemoutNbmc{}          & \textbf{\AliveCactusNumUnknownNbmc{}} & \AliveCactusGeomeanMsNbmc{} \\
  & Unb.-Ours(ric3)   & \AliveCactusNumProblemsTotal{} & \AliveCactusNumSolvedMr{}            & \textbf{\AliveCactusNumTimeoutMr{}}   & \textbf{\AliveCactusNumMemoutMr{}}   & \AliveCactusNumUnknownMr{} & \AliveCactusGeomeanMsMr{} \\
  & Unb.-Ours(k-ind.) & \AliveCactusNumProblemsTotal{} & \AliveCactusNumSolvedMk{}            & \textbf{\AliveCactusNumTimeoutMk{}}   & \textbf{\AliveCactusNumMemoutMk{}}   & \AliveCactusNumUnknownMk{} & \AliveCactusGeomeanMsMk{} \\
  & CVC5-berger         & \AliveCactusNumProblemsTotal{} & \AliveCactusNumSolvedC{}             & \AliveCactusNumTimeoutC{}             & \textbf{\AliveCactusNumMemoutC{}}    & \textbf{\AliveCactusNumUnknownC{}} & \textbf{\AliveCactusGeomeanMsC{}} \\
\midrule
\multirow{5}{*}{\shortstack[l]{InstCombine\\(zext/sext)}}
  & BMC-Ours          & \SextZextInstcombineNumProblemsTotal{} & \textbf{\SextZextInstcombineNumSolvedMbmc{}} & \SextZextInstcombineNumTimeoutMbmc{}          & \textbf{\SextZextInstcombineNumMemoutMbmc{}} & \textbf{\SextZextInstcombineNumUnknownMbmc{}} & \SextZextInstcombineGeomeanMsMbmc{} \\
  & BMC-Naive         & \SextZextInstcombineNumProblemsTotal{} & \SextZextInstcombineNumSolvedNbmc{}          & \SextZextInstcombineNumTimeoutNbmc{}          & \SextZextInstcombineNumMemoutNbmc{}          & \textbf{\SextZextInstcombineNumUnknownNbmc{}} & \SextZextInstcombineGeomeanMsNbmc{} \\
  & Unb.-Ours(ric3)   & \SextZextInstcombineNumProblemsTotal{} & \SextZextInstcombineNumSolvedMr{}            & \textbf{\SextZextInstcombineNumTimeoutMr{}}   & \textbf{\SextZextInstcombineNumMemoutMr{}}   & \SextZextInstcombineNumUnknownMr{} & \SextZextInstcombineGeomeanMsMr{} \\
  & Unb.-Ours(k-ind.) & \SextZextInstcombineNumProblemsTotal{} & \SextZextInstcombineNumSolvedMk{}            & \textbf{\SextZextInstcombineNumTimeoutMk{}}   & \textbf{\SextZextInstcombineNumMemoutMk{}}   & \SextZextInstcombineNumUnknownMk{} & \SextZextInstcombineGeomeanMsMk{} \\
  & CVC5-berger         & \SextZextInstcombineNumProblemsTotal{} & \SextZextInstcombineNumSolvedC{}             & \SextZextInstcombineNumTimeoutC{}             & \textbf{\SextZextInstcombineNumMemoutC{}}    & \SextZextInstcombineNumUnknownC{} & \textbf{\SextZextInstcombineGeomeanMsC{}} \\
\midrule
\multirow{5}{*}{ROVER}
  & BMC-Ours          & \RoverCactusNumProblemsTotal{} & \textbf{\RoverCactusNumSolvedMbmc{}} & \RoverCactusNumTimeoutMbmc{}          & \textbf{\RoverCactusNumMemoutMbmc{}} & \textbf{\RoverCactusNumUnknownMbmc{}} & \RoverCactusGeomeanMsMbmc{} \\
  & BMC-Naive         & \RoverCactusNumProblemsTotal{} & \RoverCactusNumSolvedNbmc{}          & \textbf{\RoverCactusNumTimeoutNbmc{}} & \RoverCactusNumMemoutNbmc{}          & \textbf{\RoverCactusNumUnknownNbmc{}} & \RoverCactusGeomeanMsNbmc{} \\
  & Unb.-Ours(ric3)   & \RoverCactusNumProblemsTotal{} & \RoverCactusNumSolvedMr{}            & \textbf{\RoverCactusNumTimeoutMr{}}   & \textbf{\RoverCactusNumMemoutMr{}}   & \RoverCactusNumUnknownMr{} & \RoverCactusGeomeanMsMr{} \\
  & Unb.-Ours(k-ind.) & \RoverCactusNumProblemsTotal{} & \RoverCactusNumSolvedMk{}            & \textbf{\RoverCactusNumTimeoutMk{}}   & \textbf{\RoverCactusNumMemoutMk{}}   & \RoverCactusNumUnknownMk{} & \RoverCactusGeomeanMsMk{} \\
  & CVC5-berger         & \RoverCactusNumProblemsTotal{} & \RoverCactusNumSolvedC{}             & \RoverCactusNumTimeoutC{}             & \textbf{\RoverCactusNumMemoutC{}}    & \RoverCactusNumUnknownC{} & \textbf{\RoverCactusGeomeanMsC{}} \\
\bottomrule
\end{tabular}
}
\end{table}

\subsection{Bounded Solving Eliminates Enumeration}
\label{mwbv:sec:bounded-eval}

We turn to the first question: does our $\text{PBV}_n \to \text{PBV}_1$ reduction eliminate the enumeration bottleneck?
Our reduction collapses an arbitrary number of symbolic widths
into a mono-width parameter, encoding the entire problem as one \qfbv{} query.
This makes BMC-Ours independent of the number of symbolic widths $k$,
in contrast to BMC-Naive's naive enumeration which requires $O(o^k)$ solver calls
for bound~$o$.
Our experiments (\cref{mwbv:fig:eval-compiler}) show a consistent pattern across all three datasets:
%
BMC-Ours saturates nearly every benchmark, while naive enumeration degrades
catastrophically as the number of width variables grows.
On PBV-Alive (\cref{mwbv:fig:eval-all-cactus}), where each problem has a mono-width parameter ($k=1$),
both bounded solvers are effective. Even in this setting, BMC-Ours is already an order of magnitude faster
as it avoids iterating over 64 concrete widths.

The gap widens dramatically on multi-width problems.
On InstCombine (\cref{mwbv:fig:eval-sext-zext}), with $k \geq 2$ width parameters,
BMC-Naive must issue at least $64^2 = 4096$ solver calls per problem;
it solves only \SextZextInstcombineNumSolvedNbmc{}/\SextZextInstcombineNumProblemsTotal{}.
BMC-Ours, by contrast, still nearly saturates the dataset:
\SextZextInstcombineNumSolvedMbmc{}/\SextZextInstcombineNumProblemsTotal{}.

Similarly, on ROVER (\cref{mwbv:fig:eval-all-rover}), with up to $k=12$ width variables,
naive enumeration is completely impractical: $64^{12}$ is astronomical, and
BMC-Naive solves only \RoverCactusNumSolvedNbmc{}/\RoverCactusNumProblemsTotal{}.
BMC-Ours solves \RoverCactusNumSolvedMbmc{}/\RoverCactusNumProblemsTotal{}
It is surprising that BMC-Naive has memory-outs on quite a few problems,
as we would expect it to time out before it runs out of memory.
We suspect that the raw structure of the problem may trigger the SAT solver to
consume memory more quickly than time.
%
%
The reduction, then, does eliminate the BMC enumeration bottleneck.

\subsection{Unbounded Solving Scales to Real-World Problems}
\label{mwbv:sec:unbounded-eval}

We turn to the second question: does our unbounded solver scale to real-world compiler and hardware verification problems?
While bounded solving is fast and practical,
it can only verify correctness up to a chosen bound.
Our unbounded solvers prove formulas for \emph{all} bitwidths,
and consistently outperform the state of the art CVC5-berger across every dataset.
On PBV-Alive (\cref{mwbv:fig:eval-all-cactus}), our unbounded solvers each prove \AliveCactusNumSolvedMr{} problems for all bitwidths,
versus CVC5-berger's \AliveCactusNumSolvedC{}.
On InstCombine (\cref{mwbv:fig:eval-sext-zext}), both prove \SextZextInstcombineNumSolvedMr{} versus CVC5-berger's \SextZextInstcombineNumSolvedC{}.
On ROVER (\cref{mwbv:fig:eval-all-rover}), they prove \RoverCactusNumSolvedMr{} versus CVC5-berger's \RoverCactusNumSolvedC{}.
The unsolved problems involve operations outside our linear-bitwise decidable fragment
(e.g., multiplication), for which our solvers correctly return \texttt{unknown}
rather than an unsound result.

Our automata-based unbounded solvers, then, outperform the state of the art on every benchmark.

\subsection{Diagnostic Data: Formula and Automaton Sizes}
\label{mwbv:sec:diagnostics}

\cref{mwbv:tab:eval-diagnostics} reports, as geometric means, the size of the
input $\text{PBV}_n$ formula, the mono-width $\text{PBV}_1$ \qfbv{} formula
emitted by our reduction, and the multi-width automaton state space, each as a
structural term size and an AIG (bit-level) size.
Problem size does not grow too much under the reduction, as the generated \qfbv{}
query stays a small factor larger than the input regardless of the number of
width variables. On the other hand, the automaton state space grows on the many-width ROVER
datasets, tracking the harder unbounded instances and the
residual timeouts in \cref{mwbv:tab:eval-summary}.

\begin{table}[htb]
\centering
\caption{Size diagnostics across the three datasets, reported as geometric means.
For each dataset we give the size of the input $\text{PBV}_n$ formula, the
mono-width $\text{PBV}_1$ \qfbv{} formula emitted by our reduction, and the
multi-width automaton state space, each measured both as a structural term size
(AST nodes) and as an AIG (bit-level) size. The $\text{PBV}_1$ formula stays a
small constant factor larger than the input, confirming the linear-size
reduction; the automaton state space grows on the many-width datasets, tracking
the harder unbounded instances.}
\label{mwbv:tab:eval-diagnostics}
\resizebox{\textwidth}{!}{\input{chapters/multi-width-bv/plots/output/stats/stats-table.tex}}
\end{table}

\subsection{Summary of Experimental Results}

Overall, our results (\cref{mwbv:tab:eval-summary}) answer both questions affirmatively.
While CVC5-berger consistently achieves the lowest geometric mean time on the problems it solves,
it solves fewer problems overall across all three datasets.
In our experience, it is fast on problems that work well with its rewrite rules and quantifier instantiations,
but the translation into quantified integer arithmetic seems to cause the solver to either time out or return \texttt{unknown} after lengthy solve times
on problems that are not immediately solved.
BMC-Ours saturates nearly every benchmark with a single \qfbv{} query,
evidencing that our multi-width to mono-width reduction eliminates the enumeration bottleneck.
Our automata-based unbounded solvers prove strictly more problems than CVC5-berger across all three datasets,
demonstrating that they scale to real-world compiler and hardware verification benchmarks.


\section{Discussion}
\sid{TODO: move the related work to a different place}.
\label{sec:multi:related}

The approach outlined in this chapter seems nice in theory, and to be effective in practice.
Psychologically, the completeness for the linear-bitwise fragment is quite nice.
In practice, the bounded model checker, which is an effective solver for all of parametric bitvectors,
has seen more demand amongst other researchers I have discussed these algorithms with.
This makes sense, since in practice, for production, a bounded sound and complete guarantee over a richer language appears to be more valued.

An interesting extension to this setting would be to prove bitvector appends as automata-(un)-decidable.
However, it seems implausible to me to be decidable, since, if one can encode bitvector appends, then one can encode the theory of binary strings.
If one can encode binary strings, then one can encode the theory of strings over a finite alphabet.
However, the theory of strings is poorly understood, and the best algorithms we have for deciding this are in PSPACE.
which ``feels'' like a much larger complexity class that ought to be symbolic-automata-representable.
This argument does not directly rule out the automata-based approach, since the $k$-induction solver is atempting to solve an EXPTIME problem,
which is somewhat comparable to PSPACE.

In conclusion, we have developed effective and practical algorithms for parametric bitvector solving,
and we've mechanized these in Lean. In the next chapter, we develop a mechanized solver for floating point values.
This reuses the bitblasting infrastructure, but sets it up for much more intricate reasoning.
% These two chapters complete the \emph{width} axis of the thesis, but the values
% themselves have so far remained plain machine integers. The next chapter turns
% from richer \emph{widths} to richer \emph{values}: it reuses the same verified
% bitblasting infrastructure to decide predicates over \emph{floating-point}
% numbers, whose bit-level encodings---sign, exponent, and mantissa---make them by
% far the most intricate values that compilers routinely reason about



% \paragraph{Fixed Width Bitvector Solving.}
% A wide range of techniques have been contemplated for fixed-width bitvector solving,
% using approaches such as bitblasting, word-level reasoning, and combinations thereof.
% Bitwuzla \cite{niemetz2023bitwuzla} is a state-of-the-art SMT solver for fixed-width bitvectors, that uses bitblasting
% along with word-level rewriting and preprocessing. Bitwuzla is a successor to Boolector\cite{brummayer2009boolector}.

% \paragraph{Parametric Bitvector Solving.}
% Recent work has shown renewed interest in parametric bitvector solving,
% due to applications in compiler verification and hardware optimization.
% Algorithms for parametric bitvector solving based on rewriting and reductions to nonlinear integer arithmetic \cite{bit-precise-parametric-bv} have been recently proposed
% as a practically effective solution.
% The sound and complete automata-based procedures of \cref{ch:mono-width}
% decide fragments of the PBV language in the mono-width setting ($\text{PBV}_1$), via reductions
% to automata-based decision procedures for linear integer arithmetic---an approach
% rooted in the classical result that integer arithmetic can be encoded by finite
% automata~\citep{Büchi1960}. The multi-width solver above builds on that
% automata-based approach, extending it to handle the full PBV language via our
% reduction, and improving its performance on the linear fragment via verified
% automata construction and $k$-induction.
% 
% \paragraph{Bounded Model Checking.}
% Bounded model checking (BMC) arose as a technique to prevent the exponential blowup of BDDs
% by using a SAT solver to search for counterexamples upto a given bound~\cite{biere2009bounded}.
% Our algorithm uses ideas from BMC, where we bound the mono-width parameter after our reduction, and encode the entire problem as a single \qfbv{} query.


% \section{Limitations}
% \label{sec:multi:limitations}
% 
% Two limitations bound these results, and both point beyond this chapter. The
% first is one of theory: our solvers are sound and complete for the
% \emph{linear-bitwise} fragment of the parametric theory, but nonlinear
% operations---multiplication of two symbolic-width unknowns, and the like---fall
% outside it. This is not a gap we can close by more engineering; deciding the
% nonlinear fragment requires a genuinely different tool, and we have an
% unpublished reduction to the first-order theory of the $2$-adic integers.
% The second concerns evaluation. We deliberately avoid
% preprocessing such as associativity--commutativity normalization, so as to
% measure the solving technique itself rather than a normal form that might make
% problems trivial; a production tactic would want such preprocessing, and its
% absence makes our numbers conservative. Our multi-width benchmarks are, in part,
% translated from LLVM InstCombine into parametric SMT-LIB by an agentic harness
% after we hand-generalize a representative seed set, with every generated formula
% checked against a naive enumerative solver to bound $8$ before it enters the
% dataset; this guards against mistranslation, though it does not replace a full
% correctness argument for the generalization. The mono-width results of
% \cref{ch:mono-width}, by contrast, are packaged as an artifact that has been
% independently reproduced through published artifact evaluation.

% \paragraph{From Richer Widths to Richer Values.}
% These two chapters complete the \emph{width} axis of the thesis, but the values
% themselves have so far remained plain machine integers. The next chapter turns
% from richer \emph{widths} to richer \emph{values}: it reuses the same verified
% bitblasting infrastructure to decide predicates over \emph{floating-point}
% numbers, whose bit-level encodings---sign, exponent, and mantissa---make them by
% far the most intricate values that compilers routinely reason about
% (\cref{ch:floating-point}).
